\section{Critics: what the works do not say}
\label{sec:critics}

Criticising a reading list one admires is a discipline of its own. The aim
here is constructive: for each work, what is unclear, what is assumed
silently, what a replicator would miss, and what it left for the rest of us
to trip over; then the faults the five works share.

\subsection{Grubbs (1969)}

As exposition the paper is close to faultless: the questions are organised
by use case, every test comes with a worked example, and the tables were
dependable enough to survive my Monte Carlo check half a century later. Three
criticisms stand. First, the paper's own framing invites the misuse that
became standard. It presents tests for one and for two outliers and advises
choosing the alternative in advance, but offers no operational guidance for
the analyst who does not know how many to expect; practice resolved that
ambiguity by iterating the single-outlier test, and
Section~\ref{sec:exp-grubbs} shows how completely a planted pair defeats
that habit. The proper repair arrived only with \citet{rosner1983}. Second,
the independence and normality assumptions are stated once and never
stress-tested; there is no discussion of what dependence does to the null
distribution, which is the single most relevant question for every time
series user the paper would later acquire. Third, the word outlier is
treated as self-explanatory. The paper never separates the error to be
deleted from the extreme to be studied, a distinction on which, as
\citet{campbell2013} would insist four decades later, entire archives
depend. None of this diminishes the work; all of it explains why the work
alone was never enough.

\subsection{Iglewicz and Hoaglin (1993)}

The booklet does exactly what it promises, and my experiments confirm both
halves: the robust score isolates contamination that destroys the classical
one, and the 3.5 rule behaves sensibly on the batches it was calibrated for.
The criticisms concern scope and afterlife. The recommendation is presented
with too little context for the reader to know when it stops applying: the
dependence of the score's null behaviour on sample size is treated briefly,
and dependence over time not at all, so nothing warns the engineer who
lifts the rule into a sliding window over an autocorrelated stream, which
is, on the evidence of Section~\ref{sec:exp-threshold}, precisely where it
travelled and precisely where it breaks in both directions at once. The
format compounds the problem. A slim volume in a society booklet series is
hard to obtain, so the rule is now mostly cited through secondary sources,
and folklore versions with garbled constants circulate; a method that
succeeds by being easy to apply acquires a duty to be easy to verify, and
this one is not. Finally, the booklet inherits the classical framing of
identification and rejection without engaging the accommodation alternative
its own robust machinery came from, a road not taken that
Section~\ref{sec:crit-shared} returns to.

\subsection{Shafer, Fiebrich, Arndt, Fredrickson and Hughes (2000)}

The paper's great virtue is institutional honesty: it publishes its
thresholds, its flag taxonomy, its staffing model and its failure stories,
and in doing so it defined a genre. The criticisms are those of a founding
document. The thresholds are engineering judgements, declared but not
derived, with no analysis of the false alarm and missed detection rates they
purchase; the system's error behaviour is illustrated through case studies
rather than measured through any systematic verification experiment, so a
reader cannot tell whether the battery's operating point is good, merely
that it is survivable. The spatial test's dependence on network geometry is
acknowledged but not quantified, leaving every sparser or more mountainous
network to rediscover its limits. The graded flag scale conflates severity
of suspicion with kind of fault, a design that later systems, including the
one I worked on, inherited and then had to work around with free-text
comments. And, understandably for its era, there is no code and no data, so
the replication path is institutional rather than computational: to check
the paper you must build a mesonet. The decade retrospective
\citep{fiebrich2010} repairs some of this with operational statistics, and
deserves to be read as the same work's second half.

\subsection{Campbell and colleagues (2013)}

The framework paper's clarity is its contribution, and also its limit. The
definitions, taxonomy and pipeline are crisp, but the algorithmic content is
deliberately thin: the automated checks are surveyed at the level of names,
with no guidance on parameters, ordering or interaction, so a team adopting
the framework must still invent its numerical substance, and, as this report
has argued repeatedly, the numerical substance is where quality control
actually fails. The economics of the human in the loop are asserted rather
than analysed; the paper insists on human review without estimating reviewer
load per sensor, the one quantity that decides whether the architecture
scales, and subsequent observatory practice had to push much further into
automation than the paper contemplates \citep{taylor2013}. The survey of
middleware and tools has aged fastest, as tool surveys do. And the slogan
title, effective as rhetoric, encodes a quiet bias: quantity is nothing
without quality, but a filter tuned by that maxim will always prefer the
false alarm to the miss, and the paper nowhere confronts the cost asymmetry
explicitly. These are the criticisms one makes of a successful constitution:
it needed statutes, and it left the statutes to others.

\subsection{Leigh and colleagues (2019)}

This is the methodologically strongest paper on the list, and the one whose
weaknesses are most instructive. The evaluation is built on expert labels,
and the paper is candid that the experts disagreed at the margins; every
reported sensitivity is therefore relative to a judgement standard, and the
circularity, detectors evaluated against the very human practice they are
meant to replace, is noted but not resolved, for instance by inter-labeller
agreement statistics. The study covers two sites in one hydrological
setting, so transferability, the property an agency most needs, is untested;
the class imbalance inherent in anomaly detection is handled by reporting
per-type results but not by any formal treatment of uncertainty on the rare
classes, some of which contain very few events. The feature-based branch,
statistically the most original, is also the least explainable to the
operators the framework is designed for, and the paper does not discuss how
a flag justified by a distance in feature space is to be communicated on a
trouble ticket. Finally, the framework ends at flagging; correction, which
their agency needs and later software addressed \citep{jones2022}, is out of
scope, so the decision-theoretic question, what a flag costs and what a miss
costs, remains, sixty years after \citet{anscombe1960} posed it, informal.

\subsection{Five questions I would put to the authors}
\label{sec:crit-questions}

Reading as a critic ends best in questions rather than verdicts, so, one per
work. To Grubbs: knowing that your tests would spend the next half century
running inside sliding windows over dependent data, what would you have
tabulated differently, and what warning would you have printed above the
tables? To Iglewicz and Hoaglin: is 3.5 a property of your simulations or a
property you expected of the world, and would you endorse versioning the
constant per sampling regime, as operational configurations now quietly do?
To Shafer and colleagues: across a decade of the Mesonet, what fraction of
the flags your automated battery raised turned out to be false, and what did
that number do to the attention of the person reading the morning report? To
Campbell and colleagues: how many reviewer hours per thousand sensors does
your architecture assume, and at what network size does the human in the
loop become the failure mode? To Leigh and colleagues: if a second panel of
experts relabelled your series, how much of your method ranking would
survive? None of these questions is answerable from the papers, and each of
them decides how the papers should be used, which is perhaps the sharpest
criticism a reading list can receive.

\subsection{What all five share}
\label{sec:crit-shared}

Four faults run the length of the lineage. First, threshold provenance. Every
work either derives thresholds under assumptions its users will violate, or
declares them by judgement and moves on; none proposes what operations
actually needs, a discipline of recording, justifying and re-deriving
thresholds per deployment, with the measured false alarm premium attached.
My experiments make the cost concrete from both sides: a textbook constant
that buys 1.63 false alarms per day per series, and an operational constant
that silently forfeits a third of true spikes. Second, the reviewer is the
unmodelled component. Every architecture since 2000 routes flags to a human,
yet none of the five works, nor most of their descendants, reports the
quantity that governs that human's behaviour: how many flags a reviewer
sees, how many are real, and after how many false ones they stop looking.
Alarm fatigue is a measured phenomenon in other monitoring disciplines;
environmental QC still mostly treats it as anecdote. Third, the
accommodation road not taken. The robust school offered an alternative to
identify-and-delete, estimate so that bad points cannot dominate, and the
monitoring lineage declined it for defensible institutional reasons, faults
must be named to be fixed, but the choice was never argued explicitly in any
of these works, and modern pipelines quietly reimport accommodation anyway,
through medians, MADs and robust regressions, without the vocabulary to say
so. Fourth, reproducibility tracks its era rather than leading it: exact
tables in 1969, no code in 2000, partial code and data in 2019. Only now,
with configuration-driven engines \citep{schmidt2023}, is the field
approaching the standard the subject always implied: a quality control
system whose entire judgement is an inspectable artefact.

I end the role with the criticism nearest my own work. Having spent a
semester implementing exactly this lineage, I recognise in myself the
pattern every one of these works enables: the constants I set in
Section~\ref{sec:exp-battery} felt reasonable, were unrecorded anywhere but
in code, and would have been inherited unexamined by whoever came after me.
The literature taught me the tests; none of it forced me to leave an audit
trail for the doubt itself. That, more than any new detector, seems to me
the discipline's open engineering problem.
